\section{Your Visits, Your Data, Your Robot Instructions}

\subsection{The calendar, and what it costs you}

\draftinstr{Roughly 1,000 characters introducing Figure 25. Convert the
Schedule of Activities into a calendar with durations and then total it: about
18 visits, 8 to 10 inpatient days, roughly 27 days of contact over 24 months.
State that the intervals, not the illustrative calendar dates, are what the
protocol fixes. Cite \texttt{NCICosts2024} for the practical burden framing.}

\begin{pafig}[mermaid-type: gantt]
\node[mmgoal] (a) at (0,0) {Figure 25 slot\\participant calendar};
\node[mmin] (b) at (0,-2.0) {Source:\\mermaid/fig-25-visit-timeline.md};
\draw[mmedge] (a) -- (b);
\end{pafig}
\vspace{-0.7cm}
\figcaption{Figure 25. The participant calendar from screening through the 24-month
endpoint, drawn with real durations rather than crosses in a matrix, and with the four\\
critical milestones, surgery, day 30, day 90, and the endpoint, marked to the axis.}

\draftinstr{Replace the Figure 25 slot with the mermaid-type Gantt from
\texttt{mermaid/fig-25-visit-timeline.md}. A 15 cm time axis spanning September
2026 to September 2028; five alternating section bands drawn before any bar;
3.6 mm bars with 2.2 mm clear space between rows; task labels anchored east in a
4.6 cm left gutter, never inside a bar; the four critical milestones with a
0.5pt drop line to the axis.}

\vspace{0.30em}

\noindent\begin{tabularx}{\textwidth}{L{3.3cm} C{1.4cm} C{2.0cm} C{1.9cm} Y}
\toprule
\textbf{Phase} & \textbf{Visits} & \textbf{Inpatient days} & \textbf{Clinic days} & \textbf{Notes for planning} \\
\midrule
Screening & 4 & 0 & about 4 & Imaging and tissue work; no study procedure before consent \\
Baseline and surgery & 2 & 8 to 10 & 1 & The single inpatient stay of the study \\
Early recovery & 3 & 0 & 3 & The restart advisory is a conversation, not a procedure \\
Confirmatory & 2 & 0 & 2 & Day 90 carries the two results most people want \\
Long-term & 7 & 0 & 7 & One half-day every 12 weeks for 24 months \\
\textbf{Total} & \textbf{18} & \textbf{8 to 10} & \textbf{about 17} & Roughly 27 days of contact over 24 months \\
\bottomrule
\end{tabularx}

\vspace{0.30em}

\subsection{Everything recorded about you}

\draftinstr{Roughly 1,300 characters introducing Figure 24. List every data
type: arm motion, tip force, vessel proximity across 640 sensor channels;
operative video and endoscope feed; vitals, laboratory values, CA 19-9, drain
output; the case report form. State retention in years per store, and state the
read list per store. Cite \texttt{ecfr2026part11} for the electronic-records
standard and \texttt{FDATrans2024} for the transparency principles.}

\begin{pafig}[diagrams (python)-type: pipeline with access overlay]
\node[dgtiled] (a) at (0,0) {};
\glyphdb{a}{protowhite}
\node[dglabel,anchor=north] at ([yshift=-5.4mm]a.center) {Figure 24 slot, data pipeline};
\node[dglabel,anchor=north] at (0,-1.6) {Source: diagrams-python/fig\_24\_data\_pipeline.py};
\end{pafig}
\vspace{-0.7cm}
\figcaption{Figure 24. Your data through four stages, capture, write-once hash chain, the
identified and de-identified stores, and what leaves, with the read list attached to\\
every store and no edge at all for a principal that may not read it.}

\draftinstr{Replace the Figure 24 slot with the Diagrams (Python)-type pipeline
from \texttt{diagrams-python/fig\_24\_data\_pipeline.py}. Four stage clusters
plus one access-control container; retention printed on a third label line
beneath each store tile so no leader is needed; a principal that cannot read a
store is joined by nothing at all, and the figure's foot note says that the
absence is the access statement.}

\subsection{Who can read what, and for how long}

\draftinstr{Populate this table from the \texttt{READ\_ACCESS} and
\texttt{RETENTION\_DAYS} dictionaries in
\texttt{diagrams-python/fig\_24\_data\_pipeline.py}, converting days to years.}

\vspace{0.30em}

\noindent\begin{tabularx}{\textwidth}{L{3.5cm} C{1.7cm} Y}
\toprule
\textbf{Store} & \textbf{Retention} & \textbf{Who may read it} \\
\midrule
Raw telemetry & 10 years & Site PI, Sponsor-Investigator, and you on request \\
Operative video & 10 years & Operating surgeon, Site PI, and you on request \\
Case report form & 15 years & Site PI, CRO monitor, and you on request \\
Hash chain & 15 years & Sponsor-Investigator, FDA on inspection, and you on request \\
Identified store & 15 years & Site PI, and you on request \\
De-identified store & 10 years & Sponsor-Investigator, DSMB, study statistician \\
Analysis outputs & 15 years & Everyone above, and the public at publication \\
\bottomrule
\end{tabularx}

\vspace{0.30em}

\noindent Your data is not used to train a commercial model without a separate
consent. That is a statement about what does not happen, and Figure 24 draws it
as an absent edge rather than asserting it in a caption.

\subsection{Cybersecurity, in the terms the question is asked}

\draftinstr{Roughly 900 characters answering ChatGPT concern 11 directly. The
question is "could someone hack the system", and the answer is architectural
rather than reassuring: during a procedure there is no network route out of the
building, so the attack surface is the building. State what crosses before and
after, and state what remains true if hospital connectivity fails entirely, which
is that the procedure is unaffected. Cite \texttt{FDACyber2026}.}

\subsection{The robots you will meet, and the ones you will not}

\draftinstr{Roughly 1,100 characters introducing Figure 26. The source
instruction set covers ten robot types \cite{patrobotinst}. In this protocol two
are certain, four are possible depending on recovery, and four are not used.
Knowing which is which is itself an answer to the question of what the
participant is agreeing to. State the re-scoping explicitly: the source sheets
are mixed-oncology and illustrated with raster images, and both are changed
here.}

\begin{pafig}[d2-type: card grid, full page]
\node[d2key] (a) at (0,0) {Figure 26 slot\\ten robot instruction cards};
\node[d2box] (b) at (0,-1.8) {Source:\\d2/fig-26-robot-instruction-cards.d2};
\draw[d2edge] (a) -- (b);
\end{pafig}
\vspace{-0.7cm}
\figcaption{Figure 26. Ten robot-type instruction cards re-scoped from general oncology to
this PDAC Whipple protocol, each with three numbered instructions and a re-scoping\\
note, and with the four robot types this protocol does not use marked as such.}

\draftinstr{Replace the Figure 26 slot with the full-page D2-type card grid from
\texttt{d2/fig-26-robot-instruction-cards.d2}. Two columns by five rows of
7.0 cm by 3.2 cm cards with 0.6 cm of clear space in both axes; each card a
\texttt{d2cont} with a \texttt{d2key} title band, exactly three numbered
instruction lines, and a \texttt{d2gray} foot strip; the eight-arm surgical card
drawn with a 1.2pt border and every other card at 0.8pt. No raster image.}

\subsection{Software changes during your participation}

\draftinstr{Roughly 800 characters. The software version is frozen at consent
and recorded in the version registry. Any change forces a re-consent, and until
the participant acts no study procedure proceeds under the changed version.
Declining the change ends participation and returns the participant to standard
care with no penalty. Cite \texttt{FDARealW2025} and forward to Figure 23.}
